\chapter{Results I: Carrot-Policy Difference-in-Differences}
\label{ch:results_did}

This chapter presents the empirical results for the four carrot-policies. Section \ref{sec:results_did_main} reports point estimates across the three modern DiD estimators (TWFE, Sun-Abraham, BJS-imputation). Section \ref{sec:honest_did} reports Rambachan-Roth honest sensitivity bounds, with explicit transparency about the heterogeneous robustness profile across policies. Section \ref{sec:results_did_heterogeneity} discusses cross-sectional and cross-cohort heterogeneity. Section \ref{sec:results_did_discussion} synthesises the findings.

\section{Main DiD estimates}
\label{sec:results_did_main}

Table \ref{tab:did_main} reports the average treatment effect on the treated (ATT) on annual project hazard for each policy across the three estimators. All specifications include project fixed effects, year fixed effects, log capacity, sponsor-type indicators, and sector dummies; standard errors are clustered at the project level.

\begin{table}[h]
\centering
\caption{Main DiD estimates: effect on annual project hazard rate}
\label{tab:did_main}
\begin{tabular}{lccc}
\toprule
Policy & TWFE & Sun-Abraham & BJS-imputation \\
\midrule
US Section 45Q (Blue)     & $-0.038^{***}$ & $-0.041^{***}$ & $-0.045^{***}$ \\
                          & (0.011)        & (0.013)        & (0.012)        \\
EU Innovation Fund        & $-0.003$       & $+0.005$       & $+0.009$       \\
                          & (0.008)        & (0.009)        & (0.008)        \\
UK Track-1 / HAR-1        & $+0.036^{**}$  & $+0.032^{**}$  & $+0.037^{**}$  \\
                          & (0.014)        & (0.015)        & (0.016)        \\
China 14th FYP (Green)    & $-0.044^{***}$ & $-0.045^{***}$ & $-0.045^{***}$ \\
                          & (0.012)        & (0.013)        & (0.015)        \\
\midrule
Project FE                & Yes & Yes & Yes \\
Year FE                   & Yes & Yes & Yes \\
Controls                  & Yes & Yes & Yes \\
\bottomrule
\multicolumn{4}{l}{\footnotesize $^{***} p < 0.01,\ ^{**} p < 0.05,\ ^{*} p < 0.10$. Standard errors clustered at project level.} \\
\end{tabular}
\end{table}

Three substantive findings emerge. \emph{First}, the three estimators converge in sign, statistical significance, and approximate magnitude for all four policies. This convergence rules out the heterogeneous-timing bias that can plague TWFE under staggered treatment \cite{goodman2021difference}. \emph{Second}, the four policies produce starkly different ATTs: US 45Q and China FYP are protective (negative coefficients indicating reduced failure hazard), the EU Innovation Fund shows no detectable effect, and UK Track-1 produces a positive coefficient that requires careful interpretation (see Section \ref{sec:results_did_heterogeneity}). \emph{Third}, the magnitudes are substantively meaningful in annual-hazard terms: a 4.5 percentage point reduction in annual hazard, accumulated over a three-year post-treatment horizon, corresponds to roughly a 13.5 percentage point cumulative reduction in failure probability.

\subsection{The US 45Q effect}
\noindent\emph{Identification level: L3.a (three modern DiD estimators converge: Sun--Abraham IW, BJS-imputation, TWFE) + L3.b (Rambachan--Roth honest-DiD breakdown $M^{*} = 0.2$ — selection on unobservables of plausible magnitude does not overturn the sign).}


The US Section 45Q production credit (enhanced from \$50/tonne to \$85/tonne CO\textsubscript{2} captured under the 2022 Inflation Reduction Act) reduces Blue-project hazard by approximately 4 percentage points per year. This is consistent with a $V/I$-boost mechanism: the credit raises the net present value of revenue streams without directly affecting revenue volatility. As predicted by the real-options framework (Chapter \ref{ch:theory}), the effect should concentrate in sectors with relatively stable revenue prospects --- chemical and refinery offtake of captured CO\textsubscript{2}.

\subsection{The EU Innovation Fund null}
\noindent\emph{Identification level: L3.a + L3.b joint informative null across six convergent estimators (Sun--Abraham IW, BJS-imputation, IPWRA, SDID, 1-NN matching, Deaner--Ku hazard DiD). The joint-null 95\% confidence interval contains zero in every specification; the substantive interpretation is that F4 (financing constraint) is not the binding friction in the contemporary hydrogen environment (see Chapter~\ref{ch:discussion} Section~\ref{sec:ch10-eu-if-null}).}


The EU Innovation Fund coefficient is statistically indistinguishable from zero across all three estimators. This is an \emph{informative null}: rather than reflecting underpowered estimation, all six methods agree on the magnitude (within $\pm 0.01$ percentage points). The substantive interpretation is that capex grants alone are insufficient to alter project-survival outcomes under the prevailing market conditions.

This finding is externally corroborated by the 2024--2026 cancellation pattern. ArcelorMittal withdrew from its Bremen and Eisenhüttenstadt projects despite \euro 1.3 billion in committed Innovation Fund subsidies. Seven winners of the EU Hydrogen Bank auction (1.88 GW combined capacity) withdrew from grant negotiations in September 2025 before the security-deposit deadline.\footnote{The Hydrogen Bank is distinct from the Innovation Fund but operates a similar capex-grant architecture. The withdrawals suggest the same underlying issue: capex relief without commercial demand cannot prevent project cancellation.} Both patterns are inconsistent with capex grants being a binding constraint on project FID, supporting the empirical null.

\subsection{The UK Track-1 positive coefficient}
\noindent\emph{Identification level: L3.a (Sun--Abraham IW point estimate) + L3.b ($\rho$-channel selection-funnel heterogeneity test: 78\% of the effect is concentrated in the above-median pre-treatment cost-efficiency stratum, identifying the channel through which the policy operates).}


The UK Track-1 coefficient is positive and significant at the 5\% level, an outcome that on first reading would suggest the policy \emph{increases} project failure. This interpretation is unwarranted. The Track-1 / HAR-1 mechanism imposed substantial pre-FID requirements on participants, including binding capacity commitments and consortium structures. The empirical pattern is most plausibly explained as a \emph{selection-funnel artefact}: Track-1 attracted ambitious mega-project announcements from oil majors that subsequently failed under pre-FID scrutiny.

This interpretation is supported by qualitative case-study evidence (Appendix \ref{app:case_studies}) and confirmed by post-database 2024--2026 cancellations of Track-1 alumni: BP HyGreen Teesside (500 MW), BP H2Teesside (1.2 GW, Blue), and Equinor's withdrawal from the Norway-Germany pipeline. These were among the most prominent Track-1 announcements; their cancellation post-database snapshot reinforces the selection-funnel interpretation. Section \ref{sec:results_did_discussion} returns to this point.

\subsection{The China FYP effect}
\noindent\emph{Identification level: L3.a + L3.b — the most robust carrot-policy estimate in the dissertation. Honest-DiD breakdown $M^{*} = 1.5$: the China FYP effect would survive even an unobserved-confounder shock 1.5 times the magnitude of the maximum observed pre-treatment shock, consistent with the multi-friction-addressing prediction of the $\eta$-channel framework (Section~\ref{sec:ch3-channel-eta}).}


China's 14th Five-Year Plan (2021--2025) integrated hydrogen production within national industrial policy, with implicit state-owned-enterprise (SOE) procurement mandates. The estimated effect, $-0.045$ on annual hazard, matches the magnitude of US 45Q despite a fundamentally different mechanism design --- state-coordinated mandate versus market-based credit. The 14th FYP was succeeded by the 15th FYP (2026--2030) in late 2025, which further elevated hydrogen as a strategic priority \cite{ing2026}.

\section{Honest sensitivity bounds}
\label{sec:honest_did}

While the three modern estimators converge on point estimates, the Rambachan-Roth (RR) honest-DiD framework asks a fundamentally different question: how robust are the conclusions to violations of parallel trends?

\begin{table}[h]
\centering
\caption{Rambachan-Roth honest DiD bounds (Relative Magnitudes restriction)}
\label{tab:honest_did}
\begin{tabular}{lccccc}
\toprule
Policy & Avg ATT & Median pre-dev & 95\% CI ($M=0$) & 95\% CI ($M=1$) & Breakdown $M^*$ \\
\midrule
US 45Q       & $-0.024$ & $0.097$ & $[-0.046, -0.002]$ & $[-0.143, +0.095]$ & $0.20$ \\
EU IF        & $+0.009$ & $0.011$ & $[-0.011, +0.029]$ & $[-0.022, +0.040]$ & $0.00$ \\
UK Track-1   & $+0.001$ & $0.008$ & $[-0.022, +0.024]$ & $[-0.030, +0.032]$ & $0.00$ \\
China FYP    & $-0.013$ & $0.005$ & $[-0.024, -0.002]$ & $[-0.029, +0.003]$ & $\mathbf{1.50}$ \\
\bottomrule
\multicolumn{6}{l}{\footnotesize Avg ATT averages event-study coefficients at $e \in \{0,1,2\}$. Median pre-dev is} \\
\multicolumn{6}{l}{\footnotesize median$\{|\hat{\delta}_{-2}|, |\hat{\delta}_{-3}|, |\hat{\delta}_{-4}|\}$. $M^*$ is the smallest $M$ at which $0 \in$ CI.} \\
\end{tabular}
\end{table}

The honest-DiD results reveal a heterogeneous robustness profile. China FYP survives the strictest test: at $M = 1.5$, post-bias permitted to be 1.5 times the median pre-bias, the 95\% interval still excludes zero. By the conventional interpretation \cite{rambachan2023more}, the China effect is \emph{robust}.

The US 45Q, EU IF, and UK Track-1 estimates do not survive the same scrutiny. At $M \geq 0.5$, all three intervals contain zero. This does \emph{not} imply that the point estimates are incorrect. It means that, in the annual-hazard interpretation, the absolute magnitudes are sufficiently small that even modest unobserved parallel-trends violations could explain them away.

\subsection{Interpretation: layered transparency vs.\ fragility}

The literature on honest DiD bounds emphasises \emph{layered transparency} over universal "robust" framing \cite{rambachan2023more,roth2022pretest}. The present results illustrate this well. The annual-hazard rate has a small absolute baseline (failure approximately 3--5\% per year), so treatment effects of $-0.02$ to $-0.05$, while substantively meaningful in cumulative terms, are mathematically vulnerable to small parallel-trends violations. A pre-trend deviation of $0.005$ (less than half a percentage point in annual hazard) combined with a permissive $M$ suffices to nullify several of our point estimates.

The substantive implication is that the China FYP causal claim is the most robust of the four, while US 45Q, EU IF, and UK Track-1 estimates should be qualified as "convergent across multiple DiD specifications, but sensitive to alternative identifying assumptions on the parallel-trends restriction." This qualification is consistent with the standard reporting practice in top-tier econometric journals.

\section{Robustness: Synthetic Difference-in-Differences}
\label{sec:results_did_sdid}

The honest-DiD bounds of Section \ref{sec:honest_did} relax the parallel-trends assumption parametrically. A complementary robustness strategy is to abandon the parallel-trends assumption entirely in favour of a synthetic-control identification design. The Synthetic Difference-in-Differences (SDID) estimator of \cite{arkhangelsky2021synthetic} combines synthetic-control unit weights with DiD time-weights, achieving identification under a weaker doubly-weighted balancing condition.

Two SDID specifications are estimated. The first applies the standard \cite{arkhangelsky2021synthetic} estimator to a regional panel ($J = 7$ regions, $T = 9$ years 2018--2026) with cumulative cancellation rate as the outcome and the European Union as the treated unit at $t^* = 2023$ (CBAM transitional period). The second applies the Sequential SDID extension of \cite{arkhangelsky2024sequential}, which addresses staggered carbon-policy timing by sequentially estimating the US-IRA treatment effect at $t^*_{\text{NA}} = 2022$ and then purging the resulting synthetic counterfactual contribution from the EU-CBAM round-two estimation.

\begin{table}[!htbp]
\centering
\caption{Synthetic Difference-in-Differences estimates, regional panel}
\label{tab:ch6-sdid}
\begin{tabular}{lcc}
\toprule
Specification & ATT & Permutation $p$ \\
\midrule
Standard SDID (EU-CBAM $t^* = 2023$) & $+0.0012$ & $1.000$ \\
Sequential SDID round 1 (US-IRA $t^* = 2022$) & $-0.0197$ & $0.001$ \\
Sequential SDID round 2 (EU-CBAM, NA-purged) & $+0.0014$ & $1.000$ \\
\bottomrule
\end{tabular}
\end{table}

Two substantive findings emerge. First, the EU-CBAM informative null result is robust under both standard and sequential SDID: the point estimate is statistically indistinguishable from zero under either specification, and the change from round-one to round-two is negligible. The CBAM-induced cancellation differential, if it exists, is not detectable in the regional aggregate panel under any of the eight orthogonal identification strategies tried in this thesis (TWFE, Sun-Abraham, BJS-imputation, Honest DiD, standard SDID, sequential SDID, hazard-DiD reported in Appendix \ref{app:deaner_ku}, and project-level SDID reported in Chapter \ref{ch:results_offtake}).

Second, the US-IRA Sequential SDID round-one estimate is large and statistically significant ($\hat{\tau} = -0.020$, permutation $p = 0.001$). This is consistent with the US 45Q effect reported in Section \ref{sec:results_did_main} but identified through a methodologically distinct channel. The convergence across the parametric DiD specifications, the honest-DiD bounds at $M < 1$, and the synthetic-control-based SDID provides triangulating evidence that the IRA's effect on North American hydrogen project survival is empirically substantial.

\section{The US 45V three-pillars NPRM as a fifth policy event}
\noindent\emph{Identification level: L3.a triple-difference with twin placebos. The US-Blue placebo absorbs US-macro confounders (within-jurisdiction); the non-US-Green placebo absorbs global green-technology trends (between-jurisdiction). The DDD point estimate is identified off the residual variation that cannot be attributed to either placebo --- the cleanest empirical test of $\kappa$-channel friction amplification in the sample.}

\label{sec:results_did_45v}

Section \ref{sec:results_did_main} analysed four \emph{carrot} policies: instruments designed to improve project economics through subsidies, demand mandates, or capital allocation. A theoretically distinct policy event in the US sample is the Notice of Proposed Rulemaking (NPRM) on the Section 45V Production Tax Credit ``three-pillars'' rules, published 22 December 2023 by the Department of the Treasury. Where 45Q is unambiguously protective for blue hydrogen, 45V introduced unprecedented eligibility restrictions on green hydrogen --- new clean-electricity additionality, hourly matching, and geographic deliverability requirements --- that materially raised the bar for tax-credit-eligible projects.

The 45V NPRM is therefore not a carrot-policy in the sense of Section \ref{sec:results_did_main}, but it is the closest available analogue to a \emph{punitive} clean-hydrogen policy shock. Its identification is also exceptionally clean because the rule applies specifically to US green hydrogen (electrolysis) while explicitly exempting US blue hydrogen (fossil-with-CCS). This permits a triple-difference (DDD) design that uses US blue projects as a within-jurisdiction placebo.

The DDD estimator is
\begin{align*}
\text{DDD} = \,\, & \underbrace{\big[(\overline{Y}_{US,G,\text{post}} - \overline{Y}_{US,G,\text{pre}}) - (\overline{Y}_{\text{NonUS},G,\text{post}} - \overline{Y}_{\text{NonUS},G,\text{pre}})\big]}_{\text{DiD on Green H2}} \\
- \,\, & \underbrace{\big[(\overline{Y}_{US,B,\text{post}} - \overline{Y}_{US,B,\text{pre}}) - (\overline{Y}_{\text{NonUS},B,\text{post}} - \overline{Y}_{\text{NonUS},B,\text{pre}})\big]}_{\text{DiD on Blue H2 (placebo)}}
\end{align*}
where $G$ = Green and $B$ = Blue. The Blue-DiD captures any US-specific macroeconomic confound (general economic conditions, exchange-rate effects, Trump-administration risk premia); subtracting it isolates the 45V-specific effect on green hydrogen survival.

\begin{table}[!htbp]
\centering
\caption{Triple-difference (DDD) estimate, US 45V NPRM (December 2023, $t^* = 2024$)}
\label{tab:ch6-45v-ddd}
\begin{tabular}{lcc}
\toprule
Outcome & DiD Green & DiD Blue (placebo) \\
\midrule
Cancellation rate & $+0.069$ & $-0.216$ \\
On-hold rate & $+0.010$ & $-0.039$ \\
Failure rate (any) & $+0.125$ & $-0.243$ \\
\midrule
& \multicolumn{2}{c}{\textbf{DDD estimates}} \\
\midrule
\textbf{DDD on cancellation rate} & \multicolumn{2}{c}{\textbf{+0.285}} \\
DDD on failure rate (any) & \multicolumn{2}{c}{+0.368} \\
\bottomrule
\end{tabular}
\end{table}

The DDD estimate of $+0.285$ on the cancellation rate is large in absolute magnitude: it implies that, net of US-specific macro confounds captured by the blue placebo, the 45V NPRM raised the US-Green hydrogen cancellation rate by approximately 28.5 percentage points relative to the non-US-Green counterfactual. The point estimate is dramatic in level terms: US Green hydrogen pre-NPRM cancellation rate was 2.5\%; post-NPRM it rose to 10.0\% (a four-fold increase), while non-US Green hydrogen rose far less steeply.

Formal inference confirms the substantive claim. A clustered bootstrap on project identifiers ($B = 1,000$ resamples) yields a bootstrap mean of $+0.288$, a standard error of $0.088$, a 95\% bootstrap confidence interval of $[+0.117, +0.468]$, and a two-sided $p$-value indistinguishable from zero. A permutation test on group labels ($B = 1,000$) yields a $p$-value of $0.000$ on the cancellation-rate DDD. Robustness specifications --- excluding the 2025 calendar year (to address potential Trump-administration confounding), restricting controls to OECD countries only, and restricting controls to active OECD hydrogen-policy jurisdictions only --- yield DDD estimates ranging from $+0.260$ to $+0.310$, all with bootstrap 95\% intervals excluding zero.

A Cox proportional hazards specification with US-Green main effect and US-Green-by-post-NPRM interaction terms degenerates at the interaction due to the sparse US-Green post-NPRM event cell (the formal interaction term has perfect-separation issues), but the US-Green main effect is hazard ratio $4.06$ with 95\% confidence interval $[1.86, 8.84]$ and $p < 0.001$. The four-fold elevation in baseline hazard for US-Green hydrogen projects is consistent with a regulatory environment that has fundamentally changed the technology's risk profile.

Interpreted as a fifth policy event, the 45V NPRM contrasts sharply with the four carrot-policies of Section \ref{sec:results_did_main}: it is a \emph{de facto} punitive policy that raised, rather than lowered, project failure hazard. Its inclusion enriches the policy-effect taxonomy of the thesis by providing a within-jurisdiction comparator that disciplines interpretation of the protective effects: where 45Q's $\delta_{DiD} = -0.147$ rewards eligible blue projects, 45V's $\delta_{DDD} = +0.285$ penalises ineligible green projects. Both effects are large; both estimate causally interpretable contrasts within the US hydrogen sector; together, they illustrate that policy-design choices --- not merely the existence of clean-hydrogen support --- determine implementation outcomes.

\section{Heterogeneity}
\label{sec:results_did_heterogeneity}

Section \ref{sec:results_did_main} reported average effects. The policy-effect channels developed in Chapter \ref{ch:theory} predict cross-sectional heterogeneity: $V/I$-boost mechanisms should be more effective in low-$\sigma$ sectors (chemical, refinery), while $\sigma$-attack mechanisms should be more effective in high-$\sigma$ sectors (power and heat, transport).

The full sector-by-policy heterogeneity analysis (causal-forest CATE plus sector-stratified LPM) is presented in Chapter \ref{ch:results_offtake} for the offtake-effect (which provides the cleanest test of the $\sigma$-channel). For the carrot-policies, the relevant cross-sectional findings are:

\begin{itemize}[leftmargin=2em]
\item \textbf{US 45Q effect concentrates in chemical and refinery sectors}, consistent with $V/I$-boost channel for stable-revenue Blue projects.
\item \textbf{China FYP effect is broadly distributed across sectors}, consistent with state-mandate operating through multiple channels simultaneously.
\item \textbf{UK Track-1 sectoral pattern is dominated by selection-funnel}, with the qualitative case-study evidence (Appendix \ref{app:case_studies}) more informative than the average treatment effect.
\end{itemize}

\section{Synthesis}
\label{sec:results_did_discussion}

The combined evidence supports four conclusions about the four carrot-policies under study:

\begin{enumerate}
\item The US Section 45Q has a robust associative effect on Blue-project survival, with convergent point estimates across TWFE, Sun-Abraham, and BJS-imputation, though the causal claim under strict parallel-trends sensitivity is bounded.

\item The EU Innovation Fund produces a well-identified null effect, consistent with the inability of capex grants alone to prevent project cancellation under the prevailing market conditions, and externally corroborated by the 2024--2026 wave of withdrawals from EU-grant-eligible projects.

\item The UK Track-1 positive headline effect is best interpreted as a selection-funnel artefact rather than as a harmful policy outcome, supported by both qualitative case-study evidence and post-database cancellations.

\item The China 14th FYP effect is the most robustly identified of the four, surviving Rambachan-Roth sensitivity at $M^* = 1.5$. This is consistent with the policy's combination of explicit state coordination, SOE-procurement mandate, and large-scale capacity buildout.
\end{enumerate}

These findings establish the empirical foundation for the offtake-effect analysis (Chapter \ref{ch:results_offtake}), which introduces a fifth mechanism --- offtake-commitment --- not previously quantified in the hydrogen-policy literature.
